\documentclass[12pt,letterpaper]{report}

\usepackage[utf8]{inputenc}
\usepackage[spanish]{babel}
\usepackage{graphicx}
\usepackage[left=3cm,right=3cm,top=2cm,bottom=2cm]{geometry}
\usepackage[breaklinks,colorlinks=true,linkcolor=black,citecolor=blue,urlcolor=black]{hyperref}

\makeindex

\title{Título tesis}

\renewcommand{\baselinestretch}{1.2}
\begin{document}
	\pagenumbering{roman}	
	\renewcommand{\listtablename}{Índice de tablas}
	\renewcommand{\tablename}{Tabla}
	\renewcommand{\bibname}{Referencias bibliográficas}	
%	esto es para que las listas anidadas salgan con numeros
	\renewcommand{\labelenumii}{\arabic{enumi}.\arabic{enumii}}
	\renewcommand{\labelenumiii}{\arabic{enumi}.\arabic{enumii}.\arabic{enumiii}}
	\renewcommand{\labelenumiv}{\arabic{enumi}.\arabic{enumii}.\arabic{enumiii}.\arabic{enumiv}}
	
	
	\pagestyle{empty}		
	\thispagestyle{empty} 
	
 	\begin{titlepage}

\centering

{\includegraphics[width=0.5\textwidth]{imagen/cujae}\par}

\vspace{3.5cm}

{\scshape\LARGE Sistema para el entrenamiento de operarios en la Industria Alimentaria \par}
%%{\scshape\Large SISTEMA PARA EL ENTRENAMIENTO DE OPERARIOS EN LA INDUSTRIA ALIMENTARIA \par}
\vspace{1cm}

{\bfseries\Large Trabajo de Diploma para Optar por el Título de Ingeniería en Informática\par}
\vspace{3.5cm}

{\bfseries\Large Autor:} { \Large Mónica Montoto Montané \par}
{\bfseries\Large Tutor: }{\Large Dra. Raisa Socorro Llanes \par}
\vfill

{\Large La Habana, Cuba \par}

\large{\mifecha\today}

\end{titlepage}



	%\include{dec_autoria}
	%\include{opinion_tutor}
	%\chapter*{Dedicatoria}
A mis padres, por siempre estar a mi lado, por la confianza depositada, por los empujones cuando no quería seguir, por motivarme, por hacerme una persona de bien y cada día luchar por mi futuro; a ellos, se los debo todo. Una vez me dijeron que la idea de criar hijos no es para que te acompañen cuando tú estés viejo, es para asumir la responsabilidad de criar humanos funcionales, comprometidos con la naturaleza y la sociedad. Este título también es de ustedes, porque lo sufrieron tanto como yo, y se los entrego con la promesa de que seguiré intentando ser una persona de bien, y que se sientan tan orgullosos como yo estoy de ustedes.

A mi tutora Raisa, que aceptó el reto de cuidarme y apoyarme, y lo hizo de la mejor manera. Sin dudas puedo decir, que tuve la suerte de conocer a una de las mejores profesionales que trabajan en la universidad. Gracias por enseñarme, por apoyarme y tenerme paciencia. A usted no solo la quiero, yo la admiro.
	%\chapter*{Agradecimientos}
Quiero dedicar el siguiente trabajo a todas aquellas personas que lucharon porque se hiciera realidad, y no hablo solo de la tesis, si no de todos aquellos compañeros que tuve apoyándome desde un inicio para que hoy pudiera estar aquí.
	\section*{Resumen}
La Industria Alimentaria en Cuba juega un rol muy importante dentro de la economía del país. Cuando existe una falla en su proceso productivo, no siempre están presentes los expertos para su posible solución, ya que no todos los técnicos están preparados para tomar decisiones correctas en determinadas circunstancias. En ese caso, se detiene la producción hasta que aparezca el trabajador que pueda resolver la falla ocurrida, provocando pérdidas en numerosas áreas y sectores. Es por ello que, esta industria, presenta la necesidad de capacitar a todos sus operarios ante los procesos que en ella se desarrollan.

Para brindarle una solución a esta problemática se desarrollaron, por el Instituto de Investigaciones de la Industria Alimentaria (IIIA) y la Facultad de Ingeniería Informática e Ingeniería Química de la CUJAE, un Sistema Experto para el Control de Procesos Químicos (SECPROIT) y un Generador de Bases de Conocimientos.

Este sistema experto SECPROIT, tiene como objetivo evaluar las decisiones que toman los operarios ante diferentes situaciones críticas y brinda a los jefe los resultados de las mismas y cómo se van superando sus trabajadores ante cada proceso. Actualmente, dicho sistema, presenta ciertas limitaciones que impiden su correcto funcionamiento, por lo que es necesario realizar una actualización del mismo.

Una vez logrado este objetivo se podrán resolver los problemas de la Industria Alimentaria, planteados al inicio de este resumen. A esta investigación se le atribuye un valor práctico, ya que este sistema (artefacto de salida) será implementado en las Industrias Azucareras Cubanas.

\vfill

\begin{description}
	\item[Palabras claves:]{Industria Alimentaria, proceso productivo, operario, Instituto de Investigaciones de la Industria Alimentaria (IIIA), situación crítica, Sistema Experto para el Control de Procesos Químicos (SECPROIT), Generador de Bases de Conocimientos, sistema experto, artefacto de salida.}
\end{description}



	\section*{Abstract}
The Food Industry in Cuba plays a very important role in the economy, so a problem in this sector should be avoided at all costs. It currently presents the need to train all its operators in the production processes that are developed, since it consists of many irregularities when working on them. When there is a fault, experts are not always present for its possible solution, since not all technicians are prepared to make the correct decisions in certain circumstances. In this case, the production process is stopped until an expert appears who can resolve the fault that has occurred, causing losses in many areas and thus affecting the economy.

As a result of this situation, the Food Industry Research Institute (IIIA), in conjunction with the Faculty of Computer Engineering and Chemical Engineering of CUJAE, in search of a solution, developed an Expert System for the Control of Chemical Processes ( SECPROIT) and a Knowledge Base Generator.

The SECPROIT expert system aims to evaluate the decisions made by operators in different critical situations and provides managers with their results and how their workers improve in each process. At this time, the system has certain limitations that prevent it from working properly, so it is necessary to update it. Once this objective has been achieved, the problems of the Food Industry, raised at the beginning of this summary, can be resolved.

\vfill

\begin{description}
	\item[Keywords:]{Food Industry, production process, operator, Food Industry Research Institute (IIIA), critical situation, Expert System for the Control of Chemical Processes (SECPROIT), Knowledge Base Generator, expert system.}
\end{description}
	
	\pagestyle{plain}	
	
	\tableofcontents	
	\pagebreak	
	
	\listoftables
	\pagebreak
	
	\listoffigures	
 	\clearpage 

	\pagenumbering{arabic}
	
	\phantomsection
	\addcontentsline{toc}{chapter}{Introducción}
	\pagestyle{fancy} 
	
	\chapter*{Introducción}
En la actualidad, el desarrollo tecnológico se ha convertido prácticamente en una obligación para la mayoría de los entornos laborales, industriales y sociales. En un mundo tan globalizado, que avanza y se supera con una rapidez vertiginosa, no resulta de extrañar que se exijan mejores condiciones y prestaciones. A partir de estas demandas, surge la necesidad de establecer nuevas técnicas en las grandes empresas que, de manera automática, faciliten todo tipo de actividades, desde lo laboral y comercial, hasta lo cotidiano \cite{avanceTecnologico}. 
Estas exigencias requieren nuevos desafíos, sin incluir los retos tecnológicos ya mencionados. Los avances técnicos traen consigo la necesidad de contar con factores humanos capacitados para responder ante este nuevo desarrollo científico \cite{invstCibernCuba}. 

Cuba no está ajena a esta nueva revolución tecnológica. El país está abocado a la creación y fortalecimiento de programas de posgrado que consoliden, complementen, profundicen y originen nuevos conocimientos, inmersos en áreas relacionadas con estos cambios de la era digital. De esta forma se proporcionan herramientas para afrontar retos, tomar decisiones y solucionar problemas de la sociedad cubana en general \cite{avanceTecnologico}.

El empleo de una tecnología de software propia reduce el tiempo de desarrollo y puesta marcha de la misma, minimiza el costo, aporta gran flexibilidad y eficiencia en las operaciones, mejora significativamente la trazabilidad de los procesos y facilita el trabajo a partir de diseños amigables de la interfaz hombre-máquina \cite{industriaCubana}.
Es por ello que, en Cuba, constantemente se trabaja en función de desarrollar nuevos sistemas que permitan un avance técnico significativo.
Por citar un ejemplo, la Industria Farmacéutica Cubana está compuesta por procesos tecnológicos de elevada complejidad, que necesitan cumplir con exigentes indicadores de calidad en la producción de medicamentos. A esta industria se le empleó una metodología que permitió automatizar la mayoría de sus procesos. Esta, ha significado un ahorro al país superior a tres millones y medio en moneda convertible en los últimos cinco años, un gran beneficio gracias al avance tecnológico \cite{industriaFarmceutica}.

En ocasiones, las ventajas competitivas que se obtienen con el uso de la tecnología, pueden verse afectadas por diversos riesgos, lo que provoca una posibilidad de pérdida hacia el futuro. Aunque son un hecho probabilístico sobre el que se tiene cierta incertidumbre, porque pueden ser provocados por factores internos o externos a la organización, tienen el poder de originar efectos indeseados en el cumplimiento de los objetivos de las instituciones. 
Ante esta realidad, ha surgido como necesidad la gestión de riesgos, con el propósito de lograr ejercicios de planeación y ejecución más acertados.
De esta situación no escapa la Industria Cubana y como parte de ella, la Industria Alimentaria \cite{CONVENCION2019}.

La Industria Alimentaria Cubana juega un rol muy importante dentro de la economía del país. Se caracteriza por trabajar de manera ininterrumpida, con un personal que cambia frecuentemente. Esta situación dificulta el entrenamiento de sus procesos productivos, ya que no cuenta con suficientes especialistas para la capacitación de los trabajadores. Al generarse una falla en un proceso, existe el inconveniente de una posible ausencia de expertos que encuentren una solución al problema, ya que todos los técnicos no tienen el mismo nivel de preparación para tomar decisiones ante fallas en determinadas circunstancias. En este caso, resulta entonces necesario paralizar el proceso productivo hasta localizar al experto que solucione el problema detectado, lo que genera pérdidas de recursos y un incumplimiento en el tiempo de producción \cite{gestorBases}.

El Instituto de Investigación de la Industria Alimentaria (IIIA) posee un capital humano especializado y de prestigio que a lo largo de los años, logró resultados científico-técnicos consistentes en la creación de aditivos, extensores alimenticios y la fortificación con vitaminas y minerales de alimentos seleccionados. El propósito de esta institución es lograr ventajas competitivas sostenibles en las industrias cubanas \cite{anaMailen}. 
Para darle solución al problema de la Industria Alimentaria, planteado anteriormente, este instituto, en conjunto con las facultades de Ingeniería Química e Ingeniería Informática (ambas de la CUJAE), desarrolló un Sistema Experto para el Control de Procesos Químicos (SECPROIT) y un Generador de Bases de Conocimiento.

El Generador de Bases de Conocimiento es el encargado de crear las bases de información para los diferentes procesos químicos que ocurren en una fábrica. Cada base contiene un grupo de variables, cada una con sus características y clasificaciones. De cada variable se conocen las causas que pueden provocar un estado de riesgo o alarma. De cada causa, se conocen las recomendaciones a seguir para minimizar los riesgos en el proceso productivo. A partir de estos datos, el generador elabora un conjunto de ficheros que luego serán utilizados por el SECPROIT \cite{anaMailen}.

El Sistema Experto para el Control de Procesos Químicos (SECPROIT) tiene como objetivo asesorar la toma de decisiones durante la ejecución de un proceso. A partir de los datos obtenidos por el Generador de Bases de Conocimiento, este sistema produce un conjunto de entrenamientos para los operarios de la fábrica. En cada prueba, el usuario debe escoger las variables que considere en estado de alarma, señalar sus posibles causas y seleccionar las recomendaciones a seguir. Cada prueba contiene tres etapas: variables, causas y recomendaciones. Estas evaluaciones son utilizadas para advertir los conocimientos de los trabajadores y la capacidad que posee cada uno \cite{elena}.

Actualmente el sistema SECPROIT posee ciertas limitaciones que impiden que sea implementado en las industrias:
\begin{itemize}
\item al generar o evaluar los entrenamientos, la información utilizada se extrae directamente de los ficheros obtenidos por el Generador de Bases de Conocimiento, lo que resulta una demora extra
\item solo existe un método de pregunta
\item el resultado del entrenamiento está dado por la cantidad de preguntas correctas que se respondieron, sin tener en cuenta el tiempo utilizado
\item en la última etapa, las recomendaciones, no se evalúan correctamente los resultados ni se muestra la puntuación obtenida
\item los entrenamientos se realizan de manera continua, sin posibilidad de pausa
\item las etapas aparecen de manera consecutiva, sin importar si la anterior fue aprobada o no, ya que los resultados solo influyen en la evaluación final (aparecen etapas innecesariamente)
\item los procesos se pueden evaluar una sola vez, es decir, por cada proceso existe un único entrenamiento, lo que limita la capacidad de aprendizaje del trabajador
\item los reportes del sistema no brindan toda la información requerida para conocer la capacidad del trabajador
\item las visuales del sistema son poco intuitivas, con colores oscuros y de diversos tamaños, lo que resulta poco amigable para el usuario
\end{itemize}

A partir de estas restricciones se intuyen un número de cambios necesarios a realizar. Estas modificaciones deben abarcar tanto la lógica programada, como su base de datos, sus interfaces de usuario y la estructura del sistema. Es decir, son requeridas diversas transformaciones en todas las capas internas del SECPROIT.
Es por ello que, una actualización desde el propio sistema, resulta más trabajosa que realizar un proyecto nuevo. Esto se debe a que modificar un código escrito, para adaptarlo a una nueva estructura y una nueva base de datos, resulta más complejo que elaborarlo de cero \cite{Plecka2013}.

Partiendo de este principio, resulta evidente que, la solución a esta situación problemática, es desarrollar un nuevo sistema tomando como origen el SECPROIT.
Esta actualización debe poder resolver el problema de investigación que se planteaba al principio: 
\textsl{¿Cómo lograr una completa capacitación de los operarios de la Industria Alimentaria, ante la toma de decisiones en una situación crítica?}

Por lo tanto, el objetivo general de esta investigación es: 
\begin{itemize}
\item Desarrollar un sistema capaz de adiestrar a los operarios de la Industria Alimentaria, tomando como guía el sistema SECPROIT, pero solucionando las limitaciones que este actualmente presenta.
\end{enumerate}

Para poder cumplir con el objetivo de esta investigación, deben completarse una serie de objetivos específicos, y cada uno, presenta un grupo de tareas a desarrollar:
\begin{enumerate}
\item Decidir con qué tecnologías se realizará el nuevo sistema
\begin{itemize}
\item Analizar los requisitos funcionales que debe cumplir el sistema
\item Seleccionar las tecnologías con las que se trabajará
\end{itemize}
\item Modelar la actualización del SECPROIT, tomando como base el sistema anterior
\begin{itemize}
\item Realizar diagramas auxiliares para comprender y estructurar el nuevo sistema (diagrama de caso de uso, modelo de dominio, diagrama de actividades, entre otros)
\end{itemize}
\item Implementar la nueva base de datos a partir del modelo resultante en el objetivo anterior
\begin{itemize}
\item Diseñar la nueva base de datos
\end{itemize}
\item Programar el sistema en cuestión, teniendo en cuenta las limitaciones presentes en la versión actual
\begin{itemize}
\item Implementar las funcionalidades que permitan la administración del sistema
\end{itemize}
\item Implementar las nuevas interfaces, rectificando lo señalado en el sistema anterior
\begin{itemize}
\item Realizar el diseño de las nuevas interfaces
\end{itemize}
\item Programar el entrenamiento de los operarios a partir de la información recibida por el Generador de Bases de Conocimiento
\begin{itemize}
\item Implementar la generación de preguntas asociadas a la evaluación del estado de las variables (fase de entrenamiento)
\item Implementar la generación de preguntas asociadas a la evaluación de las causas del estado de una variable (fase de entrenamiento)
\item Implementar la generación de preguntas asociadas a la evaluación de las recomendaciones del estado de una variable (fase de entrenamiento)
\item Implementar la configuración del proceso de evaluación asociado a un proceso en específico
\item Implementar la conexión con el JBOSS para obtener las respuestas correctas para las preguntas asociadas a las causas y recomendaciones
\item Implementar proceso de evaluación parcial para cada etapa
\item Implementar evaluación integral del operario en un proceso
\end{itemize}
\item Realizar pruebas al nuevo sistema resultante
\begin{itemize}
\item Diseño de las pruebas al sistema
\item Ejecución de las pruebas y experimentos al sistema
\end{itemize}
\item Desplegar la aplicación en la Facultad de Ingeniería Química, para luego implementarla en las industrias
\begin{itemize}
\item Crear un manual de usuario del nuevo sistema
\end{itemize}
\end{enumerate}

Esta investigación tiene como objeto de estudio: los sistemas de entrenamiento automatizados, los sistemas expertos y los sistemas de información basados en reglas de producción. De ellos, se centra la atención en el campo de acción que comprende el Sistema Experto para el Control de Procesos Químicos (SECPROIT).

Como artefacto de salida se obtiene la actualización de un sistema de entrenamiento (una aplicación práctica). Esta aplicación permitirá capacitar a los operarios de las industrias, específicamente a los de la Industria Alimentaria Cubana, por lo que se obtiene un valor práctico de gran importancia para la sociedad.	
	

	
\end{document}